\SourceLeaf{27}
\textsc{Gauss}'s own discussion concerns a particular division of the
$(p-1)$ residues mod.\ $p$ into three classes $S$, $S'$, $S''$, using the
representation of $p$ in the form $(mm+3nn)$. The complications involved in
determining the number $f$ in this way presumably kept \textsc{Gauss} from pursuing
the path he had taken here. Indeed, the need to extend the number system to the complex
integers $(a+b\rho)$, where the root of unity that replaces $f$ is
$\rho = \frac{-1+i\sqrt{3}}{2}$ itself, becomes especially compelling at this point.

The investigations in question were certainly carried out before 1809 and
probably followed immediately upon the developments on quadratic residues
mentioned above, which were presented in January 1808 to the Göttingen Society
of Sciences. According to notes in a diary that \textsc{Gauss} kept on his investigations
for many years from the end of March 1796\HistoricalNote{\textsuperscript{*)}}{*)}{Further details about this highly valuable diary will be given in
Vol.\ 9 of the collected works.}, he began a more detailed study of the theory of cubic
and biquadratic residues in mid-February 1807. The replacement of $f$ by
$\rho$ probably followed immediately thereafter, as the original manuscript of the
notes given above shows. It may therefore be assumed that by about the time in
question (1808) \textsc{Gauss} had definitively taken the far-reaching step of
introducing the complex integers.

\textsc{Gauss} also set out, for the complex primes $2-\rho$ and $3-\rho$, the three classes
$S$, $S'$, $S''$ reproduced above, and likewise did so for the prime divisor
$3-2\rho$ of 19. It is particularly interesting that \textsc{Gauss} gave these considerations
a geometrical form, as he also did for biquadratic residues. The complex integers
$a+b\rho$ give points in the plane that form a rhombic lattice. The lattice points
of any one system $S$ can be confined to a region that, in modern terms, can
be regarded as a domain of discontinuity of an infinite group of linear
substitutions. \textsc{Gauss} made drawings for several cases. The one for the
prime $6-\rho$ is given above. The group in question would be
\[
u' = \rho^{\nu} u + (a+b\rho),
\]
where $u$ denotes a complex variable, $\nu = 1, 2, 3$ is to be set, and $a+b\rho$ is to
range over all complex integers satisfying the congruence
\[
a+b\rho \equiv 0 \quad(\mathrm{mod.}\ 6-\rho)
\]
The regions indicated in \textsc{Gauss}'s figure refer to the group enlarged by adding
the substitution $u' = \pm u$. The arbitrary choice of the system of representatives
$S$ thus amounts to an arbitrary choice of the region boundaries.
Unfortunately, there is no indication of what \textsc{Gauss} intended by choosing
this particular shape for those boundaries.

The twelve propositions collected under [2] at the end of the notes above
follow directly upon the discussion of the numbers $a+b\rho$ in the original
manuscript. Here the system of all complex integers $a+b\rho$ is replaced
unambiguously by the system obtained by reduction mod.\ $(1+x+xx)$ of all
polynomials in $x$ with rational integer coefficients.
\par\nobreak\noindent\hfill\textsc{Fricke}.\par
\endgroup
